% --- Archon physics formalization source begin ---
% archon:physics
% archon:covers PhyXMiniProblems/problem_phyx_mini_0568.lean
% archon:source-report reports/phyx_mini/problem_phyx_mini_0568.source.json

\chapter{Physics problem phyx\_mini\_0568}
\label{ch:PhyXMiniProblems_problem_phyx_mini_0568}

\paragraph{Problem source.}
\#\# Physical scenario

A rocket ship of length \$L\_0 = 240\$ m, at rest with respect to an observer O′, travels at a constant velocity \$V = 0.6c\$ relative to an inertial observer O on the ground. A light source is placed midway in the rocket ship, at point B, between two mirrors, points A and C (see figure 17.4). The light source emits two flashes, one in the direction of the nose of the ship (point C) and one in the direction of the tail (point A). Subsequently, the two flashes return to the source where they are absorbed.

\#\# Question

Determine the time that the flashes take to reach the mirrors as measured by an observer O′ on the rocket ship.

\#\# Answer choices

- A: A: 0.55\textbackslash{}mu\textbackslash{}s

- B: B: 0.4\textbackslash{}mu\textbackslash{}s

- C: C: 0.95\textbackslash{}mu\textbackslash{}s

- D: D: 10.2\textbackslash{}mu\textbackslash{}s

\#\# Figure caption (auxiliary; use the image as primary evidence)

The image is a diagram featuring a coordinate system, a rocket, and labeled points. Here's a detailed description:

- **Coordinate System**: 
  - There are two coordinate systems: (x, y) and (x', y').
  - Both have horizontal x-axes, aligned with each other as x ≡ x'.
  - The y-axis is vertical and intersects with the y'-axis.
  - Points O and O' represent the origins of these coordinate systems.

- **Rocket**:
  - Positioned horizontally along the x ≡ x' axis.
  - It is depicted with three circular windows and appears to be moving to the right in the positive x-direction.

- **Velocity**:
  - A vector labeled \textbackslash{}( V \textbackslash{}) is shown, indicating movement along the x-axis.

- **Labeled Points and Distances**:
  - Points A, B, and C are marked near the rocket.
  - Two red arrows indicate distances, with "L₀/2" shown between A and B, and between B and C.
  - A longer arrow labeled "L₀" extends from the front to the back of the rocket, along its length.

This diagram likely relates to concepts of motion or relativity, involving a reference frame, possibly illustrating length contraction or similar phenomena.

\#\# Dataset metadata

Relativity; Physical Model Grounding Reasoning, Multi-Formula Reasoning

\paragraph{Recorded answer/context.}
B: B: 0.4\textbackslash{}mu\textbackslash{}s

\paragraph{Figure/image path.}
/root/proposal\_for\_physic/science-mango/phyx\_mini\_run/phyx\_data/test\_image/568.png

\paragraph{Formalization target.}
create a compiling Lean file with sorry bodies at `PhyXMiniProblems/problem\_phyx\_mini\_0568.lean`. The Lean declarations must preserve the physical quantities, dimensions or dimensional roles, figure labels, governing-law hypotheses, and final relation expressed by this problem.
Use Mathlib/Physlib names found through LeanExplore where available. If a physics API is missing, introduce faithful local abstractions rather than scalar placeholder aliases.

\begin{theorem}[Physics formalization target]
\label{thm:physics:phyx_mini_0568:target}
\lean{PhyXMiniProblems.ProblemPhyXMini0568.problem_phyx_mini_0568}
\uses{def:physics:phyx-mini-0568:phyxminiproblems-problemphyxmini0568-timeinmicroseconds, def:physics:phyx-mini-0568:phyxminiproblems-problemphyxmini0568-vacuumlightspeedinmeterspermicrosecond, def:physics:phyx-mini-0568:phyxminiproblems-problemphyxmini0568-flashlabel, def:physics:phyx-mini-0568:phyxminiproblems-problemphyxmini0568-rocketlighttravelsetup, def:physics:phyx-mini-0568:phyxminiproblems-problemphyxmini0568-matchesrocketlightscenario, def:physics:phyx-mini-0568:phyxminiproblems-problemphyxmini0568-matchessuppliedrocketlightfigure, def:physics:phyx-mini-0568:phyxminiproblems-problemphyxmini0568-matchesproblemreadouts, def:physics:phyx-mini-0568:phyxminiproblems-problemphyxmini0568-satisfiesmidpointgeometry, def:physics:phyx-mini-0568:phyxminiproblems-problemphyxmini0568-hasphysicalrocketlightparameters, def:physics:phyx-mini-0568:phyxminiproblems-problemphyxmini0568-satisfiesrocketframelightpropagation, def:physics:phyx-mini-0568:phyxminiproblems-problemphyxmini0568-answerchoice, def:physics:phyx-mini-0568:phyxminiproblems-problemphyxmini0568-matchesanswerchoice}
The assigned autoformalize agent should translate this physics problem into Lean declarations in the covered file, with theorem and lemma proof bodies written as `by sorry`.
\end{theorem}
\begin{proof}
This is an autoformalization task, not a proof task. Produce statements that can later be proved by the physics prover without weakening the physical model.
\end{proof}

% --- Archon declaration topology begin ---
\section{Declaration topology}
\begin{definition}[Length Quantity]
  \label{def:physics:phyx-mini-0568:phyxminiproblems-problemphyxmini0568-lengthquantity}
  \lean{PhyXMiniProblems.ProblemPhyXMini0568.LengthQuantity}
  A nonnegative physical length, independent of the unit used to read it
\end{definition}
\begin{proof}
  This is the declared model definition of Length Quantity.
\end{proof}
\begin{definition}[Time Quantity]
  \label{def:physics:phyx-mini-0568:phyxminiproblems-problemphyxmini0568-timequantity}
  \lean{PhyXMiniProblems.ProblemPhyXMini0568.TimeQuantity}
  A coordinate-time reading or elapsed duration
\end{definition}
\begin{proof}
  This is the declared model definition of Time Quantity.
\end{proof}
\begin{definition}[Speed Quantity]
  \label{def:physics:phyx-mini-0568:phyxminiproblems-problemphyxmini0568-speedquantity}
  \lean{PhyXMiniProblems.ProblemPhyXMini0568.SpeedQuantity}
  A signed physical speed along the common longitudinal axis
\end{definition}
\begin{proof}
  This is the declared model definition of Speed Quantity.
\end{proof}
\begin{definition}[length Readout]
  \label{def:physics:phyx-mini-0568:phyxminiproblems-problemphyxmini0568-lengthreadout}
  \lean{PhyXMiniProblems.ProblemPhyXMini0568.lengthReadout}
  \uses{def:physics:phyx-mini-0568:phyxminiproblems-problemphyxmini0568-lengthquantity}
  Read a physical length in a selected length unit
\end{definition}
\begin{proof}
  This is the declared model definition of length Readout.
\end{proof}
\begin{definition}[time Readout]
  \label{def:physics:phyx-mini-0568:phyxminiproblems-problemphyxmini0568-timereadout}
  \lean{PhyXMiniProblems.ProblemPhyXMini0568.timeReadout}
  \uses{def:physics:phyx-mini-0568:phyxminiproblems-problemphyxmini0568-timequantity}
  Read a time quantity in a selected time unit
\end{definition}
\begin{proof}
  This is the declared model definition of time Readout.
\end{proof}
\begin{definition}[speed Readout]
  \label{def:physics:phyx-mini-0568:phyxminiproblems-problemphyxmini0568-speedreadout}
  \lean{PhyXMiniProblems.ProblemPhyXMini0568.speedReadout}
  \uses{def:physics:phyx-mini-0568:phyxminiproblems-problemphyxmini0568-speedquantity}
  Read a speed in selected length and time units
\end{definition}
\begin{proof}
  This is the declared model definition of speed Readout.
\end{proof}
\begin{definition}[length In Meters]
  \label{def:physics:phyx-mini-0568:phyxminiproblems-problemphyxmini0568-lengthinmeters}
  \lean{PhyXMiniProblems.ProblemPhyXMini0568.lengthInMeters}
  \uses{def:physics:phyx-mini-0568:phyxminiproblems-problemphyxmini0568-lengthquantity, def:physics:phyx-mini-0568:phyxminiproblems-problemphyxmini0568-lengthreadout}
  Metre readout used for the stated proper length
\end{definition}
\begin{proof}
  This is the declared model definition of length In Meters.
\end{proof}
\begin{definition}[time In Microseconds]
  \label{def:physics:phyx-mini-0568:phyxminiproblems-problemphyxmini0568-timeinmicroseconds}
  \lean{PhyXMiniProblems.ProblemPhyXMini0568.timeInMicroseconds}
  \uses{def:physics:phyx-mini-0568:phyxminiproblems-problemphyxmini0568-timequantity, def:physics:phyx-mini-0568:phyxminiproblems-problemphyxmini0568-timereadout}
  Microsecond readout used by the four answer choices
\end{definition}
\begin{proof}
  This is the declared model definition of time In Microseconds.
\end{proof}
\begin{definition}[speed In Meters Per Microsecond]
  \label{def:physics:phyx-mini-0568:phyxminiproblems-problemphyxmini0568-speedinmeterspermicrosecond}
  \lean{PhyXMiniProblems.ProblemPhyXMini0568.speedInMetersPerMicrosecond}
  \uses{def:physics:phyx-mini-0568:phyxminiproblems-problemphyxmini0568-speedquantity, def:physics:phyx-mini-0568:phyxminiproblems-problemphyxmini0568-speedreadout}
  Metre-per-microsecond readout of a signed speed
\end{definition}
\begin{proof}
  This is the declared model definition of speed In Meters Per Microsecond.
\end{proof}
\begin{definition}[vacuum Light Speed In Meters Per Microsecond]
  \label{def:physics:phyx-mini-0568:phyxminiproblems-problemphyxmini0568-vacuumlightspeedinmeterspermicrosecond}
  \lean{PhyXMiniProblems.ProblemPhyXMini0568.vacuumLightSpeedInMetersPerMicrosecond}
  \uses{def:physics:phyx-mini-0568:phyxminiproblems-problemphyxmini0568-speedinmeterspermicrosecond}
  Physlib's exact vacuum light speed, read in metres per microsecond
\end{definition}
\begin{proof}
  This is the declared model definition of vacuum Light Speed In Meters Per Microsecond.
\end{proof}
\begin{definition}[Inertial Frame Label]
  \label{def:physics:phyx-mini-0568:phyxminiproblems-problemphyxmini0568-inertialframelabel}
  \lean{PhyXMiniProblems.ProblemPhyXMini0568.InertialFrameLabel}
  The ground frame O and the rocket rest frame O'
\end{definition}
\begin{proof}
  This is the declared model definition of Inertial Frame Label.
\end{proof}
\begin{definition}[Observer Label]
  \label{def:physics:phyx-mini-0568:phyxminiproblems-problemphyxmini0568-observerlabel}
  \lean{PhyXMiniProblems.ProblemPhyXMini0568.ObserverLabel}
  The two observers named in the problem and shown in the figure
\end{definition}
\begin{proof}
  This is the declared model definition of Observer Label.
\end{proof}
\begin{definition}[Rocket Point Label]
  \label{def:physics:phyx-mini-0568:phyxminiproblems-problemphyxmini0568-rocketpointlabel}
  \lean{PhyXMiniProblems.ProblemPhyXMini0568.RocketPointLabel}
  The labeled rocket points: tail mirror A, source B, and nose mirror C
\end{definition}
\begin{proof}
  This is the declared model definition of Rocket Point Label.
\end{proof}
\begin{definition}[Flash Label]
  \label{def:physics:phyx-mini-0568:phyxminiproblems-problemphyxmini0568-flashlabel}
  \lean{PhyXMiniProblems.ProblemPhyXMini0568.FlashLabel}
  The two flashes are distinguished by which mirror they approach first
\end{definition}
\begin{proof}
  This is the declared model definition of Flash Label.
\end{proof}
\begin{definition}[Flash Event]
  \label{def:physics:phyx-mini-0568:phyxminiproblems-problemphyxmini0568-flashevent}
  \lean{PhyXMiniProblems.ProblemPhyXMini0568.FlashEvent}
  The three successive events tracked for each flash
\end{definition}
\begin{proof}
  This is the declared model definition of Flash Event.
\end{proof}
\begin{definition}[Horizontal Direction]
  \label{def:physics:phyx-mini-0568:phyxminiproblems-problemphyxmini0568-horizontaldirection}
  \lean{PhyXMiniProblems.ProblemPhyXMini0568.HorizontalDirection}
  Directions along the common horizontal axis
\end{definition}
\begin{proof}
  This is the declared model definition of Horizontal Direction.
\end{proof}
\begin{definition}[Figure Axis]
  \label{def:physics:phyx-mini-0568:phyxminiproblems-problemphyxmini0568-figureaxis}
  \lean{PhyXMiniProblems.ProblemPhyXMini0568.FigureAxis}
  Coordinate axes explicitly drawn for the two frames
\end{definition}
\begin{proof}
  This is the declared model definition of Figure Axis.
\end{proof}
\begin{definition}[Figure Text Label]
  \label{def:physics:phyx-mini-0568:phyxminiproblems-problemphyxmini0568-figuretextlabel}
  \lean{PhyXMiniProblems.ProblemPhyXMini0568.FigureTextLabel}
  Literal textual labels visible in figure 17.4
\end{definition}
\begin{proof}
  This is the declared model definition of Figure Text Label.
\end{proof}
\begin{definition}[Figure Segment]
  \label{def:physics:phyx-mini-0568:phyxminiproblems-problemphyxmini0568-figuresegment}
  \lean{PhyXMiniProblems.ProblemPhyXMini0568.FigureSegment}
  The three longitudinal segments whose lengths are annotated in the figure
\end{definition}
\begin{proof}
  This is the declared model definition of Figure Segment.
\end{proof}
\begin{definition}[Figure Distance Label]
  \label{def:physics:phyx-mini-0568:phyxminiproblems-problemphyxmini0568-figuredistancelabel}
  \lean{PhyXMiniProblems.ProblemPhyXMini0568.FigureDistanceLabel}
  Symbolic length annotations printed beside the red dimension arrows
\end{definition}
\begin{proof}
  This is the declared model definition of Figure Distance Label.
\end{proof}
\begin{definition}[Rocket Light Figure]
  \label{def:physics:phyx-mini-0568:phyxminiproblems-problemphyxmini0568-rocketlightfigure}
  \lean{PhyXMiniProblems.ProblemPhyXMini0568.RocketLightFigure}
  \uses{def:physics:phyx-mini-0568:phyxminiproblems-problemphyxmini0568-horizontaldirection, def:physics:phyx-mini-0568:phyxminiproblems-problemphyxmini0568-figureaxis, def:physics:phyx-mini-0568:phyxminiproblems-problemphyxmini0568-figuretextlabel, def:physics:phyx-mini-0568:phyxminiproblems-problemphyxmini0568-figuresegment, def:physics:phyx-mini-0568:phyxminiproblems-problemphyxmini0568-figuredistancelabel}
  Primary-image information: axes, origins, point ordering, velocity arrow, and the L₀ and L₀/2 annotations. It contains no travel-time readout
\end{definition}
\begin{proof}
  This is the declared model definition of Rocket Light Figure.
\end{proof}
\begin{definition}[Rocket Light Travel Setup]
  \label{def:physics:phyx-mini-0568:phyxminiproblems-problemphyxmini0568-rocketlighttravelsetup}
  \lean{PhyXMiniProblems.ProblemPhyXMini0568.RocketLightTravelSetup}
  \uses{def:physics:phyx-mini-0568:phyxminiproblems-problemphyxmini0568-lengthquantity, def:physics:phyx-mini-0568:phyxminiproblems-problemphyxmini0568-timequantity, def:physics:phyx-mini-0568:phyxminiproblems-problemphyxmini0568-speedquantity, def:physics:phyx-mini-0568:phyxminiproblems-problemphyxmini0568-inertialframelabel, def:physics:phyx-mini-0568:phyxminiproblems-problemphyxmini0568-observerlabel, def:physics:phyx-mini-0568:phyxminiproblems-problemphyxmini0568-rocketpointlabel, def:physics:phyx-mini-0568:phyxminiproblems-problemphyxmini0568-flashlabel, def:physics:phyx-mini-0568:phyxminiproblems-problemphyxmini0568-flashevent, def:physics:phyx-mini-0568:phyxminiproblems-problemphyxmini0568-horizontaldirection, def:physics:phyx-mini-0568:phyxminiproblems-problemphyxmini0568-rocketlightfigure}
  The independent physical data of the experiment. The outbound and return durations are unknown physical quantities. They are constrained only by the general light-propagation laws below, rather than defined from an answer
\end{definition}
\begin{proof}
  This is the declared model definition of Rocket Light Travel Setup.
\end{proof}
\begin{definition}[Matches Rocket Light Scenario]
  \label{def:physics:phyx-mini-0568:phyxminiproblems-problemphyxmini0568-matchesrocketlightscenario}
  \lean{PhyXMiniProblems.ProblemPhyXMini0568.MatchesRocketLightScenario}
  \uses{def:physics:phyx-mini-0568:phyxminiproblems-problemphyxmini0568-timereadout, def:physics:phyx-mini-0568:phyxminiproblems-problemphyxmini0568-rocketlighttravelsetup}
  Frame assignments, physical roles, flash directions, and event incidences from the prose. Simultaneous emission is stated as an event fact in O'
\end{definition}
\begin{proof}
  This is the declared model definition of Matches Rocket Light Scenario.
\end{proof}
\begin{definition}[Matches Supplied Rocket Light Figure]
  \label{def:physics:phyx-mini-0568:phyxminiproblems-problemphyxmini0568-matchessuppliedrocketlightfigure}
  \lean{PhyXMiniProblems.ProblemPhyXMini0568.MatchesSuppliedRocketLightFigure}
  \uses{def:physics:phyx-mini-0568:phyxminiproblems-problemphyxmini0568-rocketlighttravelsetup}
  Exact qualitative and symbolic information transcribed from the image
\end{definition}
\begin{proof}
  This is the declared model definition of Matches Supplied Rocket Light Figure.
\end{proof}
\begin{definition}[Matches Problem Readouts]
  \label{def:physics:phyx-mini-0568:phyxminiproblems-problemphyxmini0568-matchesproblemreadouts}
  \lean{PhyXMiniProblems.ProblemPhyXMini0568.MatchesProblemReadouts}
  \uses{def:physics:phyx-mini-0568:phyxminiproblems-problemphyxmini0568-speedreadout, def:physics:phyx-mini-0568:phyxminiproblems-problemphyxmini0568-lengthinmeters, def:physics:phyx-mini-0568:phyxminiproblems-problemphyxmini0568-rocketlighttravelsetup}
  The numerical readouts stated in the problem: L₀ = 240 m and V = 0.6c = (3/5)c. The speed relation is required in every compatible pair of length and time units; it does not determine either requested time
\end{definition}
\begin{proof}
  This is the declared model definition of Matches Problem Readouts.
\end{proof}
\begin{definition}[Satisfies Midpoint Geometry]
  \label{def:physics:phyx-mini-0568:phyxminiproblems-problemphyxmini0568-satisfiesmidpointgeometry}
  \lean{PhyXMiniProblems.ProblemPhyXMini0568.SatisfiesMidpointGeometry}
  \uses{def:physics:phyx-mini-0568:phyxminiproblems-problemphyxmini0568-lengthreadout, def:physics:phyx-mini-0568:phyxminiproblems-problemphyxmini0568-flashlabel, def:physics:phyx-mini-0568:phyxminiproblems-problemphyxmini0568-rocketlighttravelsetup}
  The source is midway between the mirrors in the rocket rest frame. Thus both physical source-to-mirror distances are half the proper rocket length. This is figure/scenario geometry, not a travel-time conclusion
\end{definition}
\begin{proof}
  This is the declared model definition of Satisfies Midpoint Geometry.
\end{proof}
\begin{definition}[Has Physical Rocket Light Parameters]
  \label{def:physics:phyx-mini-0568:phyxminiproblems-problemphyxmini0568-hasphysicalrocketlightparameters}
  \lean{PhyXMiniProblems.ProblemPhyXMini0568.HasPhysicalRocketLightParameters}
  \uses{def:physics:phyx-mini-0568:phyxminiproblems-problemphyxmini0568-lengthinmeters, def:physics:phyx-mini-0568:phyxminiproblems-problemphyxmini0568-timeinmicroseconds, def:physics:phyx-mini-0568:phyxminiproblems-problemphyxmini0568-speedinmeterspermicrosecond, def:physics:phyx-mini-0568:phyxminiproblems-problemphyxmini0568-rocketlighttravelsetup}
  Positivity, subluminality, and chronological order of both flash paths
\end{definition}
\begin{proof}
  This is the declared model definition of Has Physical Rocket Light Parameters.
\end{proof}
\begin{definition}[Satisfies Rocket Frame Light Propagation]
  \label{def:physics:phyx-mini-0568:phyxminiproblems-problemphyxmini0568-satisfiesrocketframelightpropagation}
  \lean{PhyXMiniProblems.ProblemPhyXMini0568.SatisfiesRocketFrameLightPropagation}
  \uses{def:physics:phyx-mini-0568:phyxminiproblems-problemphyxmini0568-lengthreadout, def:physics:phyx-mini-0568:phyxminiproblems-problemphyxmini0568-timereadout, def:physics:phyx-mini-0568:phyxminiproblems-problemphyxmini0568-speedreadout, def:physics:phyx-mini-0568:phyxminiproblems-problemphyxmini0568-flashlabel, def:physics:phyx-mini-0568:phyxminiproblems-problemphyxmini0568-rocketlighttravelsetup}
  Elapsed times are event-time differences in O'. On both the outgoing and returning legs, the source--mirror distance equals the invariant vacuum light speed times the elapsed rocket-frame time. These general laws mention neither 0.4 μs nor any answer label
\end{definition}
\begin{proof}
  This is the declared model definition of Satisfies Rocket Frame Light Propagation.
\end{proof}
\begin{definition}[Answer Choice]
  \label{def:physics:phyx-mini-0568:phyxminiproblems-problemphyxmini0568-answerchoice}
  \lean{PhyXMiniProblems.ProblemPhyXMini0568.AnswerChoice}
  Labels of the four microsecond-valued answer choices
\end{definition}
\begin{proof}
  This is the declared model definition of Answer Choice.
\end{proof}
\begin{definition}[displayed Time In Microseconds]
  \label{def:physics:phyx-mini-0568:phyxminiproblems-problemphyxmini0568-displayedtimeinmicroseconds}
  \lean{PhyXMiniProblems.ProblemPhyXMini0568.displayedTimeInMicroseconds}
  \uses{def:physics:phyx-mini-0568:phyxminiproblems-problemphyxmini0568-answerchoice}
  Time in microseconds displayed beside each answer label
\end{definition}
\begin{proof}
  This is the declared model definition of displayed Time In Microseconds.
\end{proof}
\begin{definition}[recorded Dataset Answer]
  \label{def:physics:phyx-mini-0568:phyxminiproblems-problemphyxmini0568-recordeddatasetanswer}
  \lean{PhyXMiniProblems.ProblemPhyXMini0568.recordedDatasetAnswer}
  \uses{def:physics:phyx-mini-0568:phyxminiproblems-problemphyxmini0568-answerchoice}
  The answer label recorded by the source dataset, retained only as metadata
\end{definition}
\begin{proof}
  This is the declared model definition of recorded Dataset Answer.
\end{proof}
\begin{definition}[Agrees With Displayed Microsecond Value]
  \label{def:physics:phyx-mini-0568:phyxminiproblems-problemphyxmini0568-agreeswithdisplayedmicrosecondvalue}
  \lean{PhyXMiniProblems.ProblemPhyXMini0568.AgreesWithDisplayedMicrosecondValue}
  Agreement with a displayed microsecond value to within 0.005 μs, enough to interpret the finite-precision choices without replacing the exact value
\end{definition}
\begin{proof}
  This is the declared model definition of Agrees With Displayed Microsecond Value.
\end{proof}
\begin{definition}[Matches Answer Choice]
  \label{def:physics:phyx-mini-0568:phyxminiproblems-problemphyxmini0568-matchesanswerchoice}
  \lean{PhyXMiniProblems.ProblemPhyXMini0568.MatchesAnswerChoice}
  \uses{def:physics:phyx-mini-0568:phyxminiproblems-problemphyxmini0568-timeinmicroseconds, def:physics:phyx-mini-0568:phyxminiproblems-problemphyxmini0568-flashlabel, def:physics:phyx-mini-0568:phyxminiproblems-problemphyxmini0568-rocketlighttravelsetup, def:physics:phyx-mini-0568:phyxminiproblems-problemphyxmini0568-answerchoice, def:physics:phyx-mini-0568:phyxminiproblems-problemphyxmini0568-displayedtimeinmicroseconds, def:physics:phyx-mini-0568:phyxminiproblems-problemphyxmini0568-agreeswithdisplayedmicrosecondvalue}
  Both rocket-frame outbound durations agree with one displayed choice
\end{definition}
\begin{proof}
  This is the declared model definition of Matches Answer Choice.
\end{proof}
\begin{lemma}[outbound Travel Time Formula]
  \label{lem:physics:phyx-mini-0568:phyxminiproblems-problemphyxmini0568-outboundtraveltimeformula}
  \lean{PhyXMiniProblems.ProblemPhyXMini0568.outboundTravelTimeFormula}
  \uses{def:physics:phyx-mini-0568:phyxminiproblems-problemphyxmini0568-lengthinmeters, def:physics:phyx-mini-0568:phyxminiproblems-problemphyxmini0568-timeinmicroseconds, def:physics:phyx-mini-0568:phyxminiproblems-problemphyxmini0568-vacuumlightspeedinmeterspermicrosecond, def:physics:phyx-mini-0568:phyxminiproblems-problemphyxmini0568-flashlabel, def:physics:phyx-mini-0568:phyxminiproblems-problemphyxmini0568-rocketlighttravelsetup, def:physics:phyx-mini-0568:phyxminiproblems-problemphyxmini0568-satisfiesmidpointgeometry, def:physics:phyx-mini-0568:phyxminiproblems-problemphyxmini0568-hasphysicalrocketlightparameters, def:physics:phyx-mini-0568:phyxminiproblems-problemphyxmini0568-satisfiesrocketframelightpropagation}
  The midpoint geometry and invariant light-speed law give (L₀/2)/c for each outbound leg. This proof-route lemma has no separate
\end{lemma}
\begin{proof}
  The conclusion follows from the governing relations and figure readouts named in the statement of outbound Travel Time Formula.
\end{proof}
% --- Archon declaration topology end ---
% --- Archon physics formalization source end ---
